\chapter{Methodology}
\label{ch:methodology}

%% ── 4.1 ─────────────────────────────────────────────────────────────────────
\section{Design Principles and System Overview}
\label{sec:design-overview}

\subsection{Guiding Design Principles}
\label{subsec:design-principles}

% Four principles that governed every design decision:
% 1. CONTINUITY: reuse same hardware + sensor suite from prior project
% 2. COMPATIBILITY: all FL messages must fit LoRaWAN payload limits (51 bytes at DR0)
% 3. MINIMALITY: model + runtime must fit in 32KB SRAM alongside firmware
% 4. PRIVACY BY DESIGN: raw sensor data never leaves the device

\subsection{Three-Layer System Architecture}
\label{subsec:system-overview}

% THREE-LAYER ARCHITECTURE (from PAPER_README.md Section 1.4):
% Layer 1 — Model Preparation (offline): train_model.py → model.h
% Layer 2 — Node Intelligence (online): FederatedTinyML.ino → sense, infer, train, send
% Layer 3 — Federated Orchestration: fl_server.py → aggregate, redistribute

% FIGURE: Full system architecture diagram showing all three layers
% \begin{figure}[H]
%   \centering
%   \includegraphics[width=\textwidth]{figures/system_architecture.pdf}
%   \caption{Three-layer system architecture.}
%   \label{fig:system-arch}
% \end{figure}

%% ── 4.2 ─────────────────────────────────────────────────────────────────────
\section{Hardware Platform and Deployment}
\label{sec:hardware}

\subsection{Arduino MKR WAN 1310}
% SAMD21 Cortex-M0+ at 48 MHz
% 32 KB SRAM, 256 KB Flash
% Murata CMWX1ZZABZ-078 LoRa module (SX1276)
% EU868 band, Class A, OTAA
% Same board used in prior project — ensures direct comparability

\subsection{Sensor Suite}
% Same sensors as prior project:
\begin{table}[H]
  \centering
  \caption{Sensor specifications (identical to prior project deployment).}
  \label{tab:sensors}
  \begin{tabular}{llll}
    \toprule
    \textbf{Sensor} & \textbf{Measurement} & \textbf{Range} & \textbf{Interface} \\
    \midrule
    BME280  & Pressure    & 300--1100\,hPa       & I\textsuperscript{2}C (0x77) \\
    SCD4x   & CO\textsubscript{2}  & 400--5000\,ppm & I\textsuperscript{2}C \\
    SCD4x   & Temperature & $-10$ to $+60$\,°C   & I\textsuperscript{2}C \\
    SCD4x   & Humidity    & 0--100\,\%RH          & I\textsuperscript{2}C \\
    SPS30   & PM2.5       & 0--1000\,µg/m³        & I\textsuperscript{2}C \\
    \bottomrule
  \end{tabular}
\end{table}

\subsection{Deployment Topology}
% 6 nodes: ED0–ED5 across indoor office rooms
% Locations: stairwell, corridor, kitchen, office rooms, etc.
% Single gateway
% FIGURE: Floor plan with node locations and distances from gateway

%% ── 4.3 ─────────────────────────────────────────────────────────────────────
\section{Dataset}
\label{sec:dataset}

\subsection{Measurement Campaign}
% Dates: September 2024 — May 2025 (8 months)
% 6 Arduino MKR WAN 1310 devices (ED0–ED5)
% Rooms: stairwell, corridor, kitchen, offices, etc.
% Raw TTN export: 2,313,903 rows, 81 columns, 1.95 GB
% Aggregated ML-ready: 1,715,869 rows, 20 columns, 298 MB

\subsection{Data Quality and Cleaning}
% Deterministic anomaly removal (33 rows, 3 patterns):
%   Pattern A: co2=21547, humidity=156.65, temp=174.90, pressure=3.21, pm25=33.93
%   Pattern B: co2=16724, humidity=210.53, temp=110.76, pressure=317.45, pm25=125.57
%   Pattern C: co2=0, humidity=0, temp=0, pressure=508.90, pm25=0
% Pressure correction: stored_value × 3.125 → real hPa
% Null handling: 1,427 missing SNR, 19 missing f_count → dropped
% Time-based sorting for temporal consistency

\begin{table}[H]
  \centering
  \caption{Data preprocessing pipeline: row counts at each stage.}
  \label{tab:data-cleaning}
  \begin{tabular}{lr}
    \toprule
    \textbf{Stage} & \textbf{Rows} \\
    \midrule
    Raw TTN export                    & 2,313,903 \\
    After aggregation                 & 1,715,869 \\
    After anomaly removal ($-33$)     & 1,715,836 \\
    After null SNR/f\_count removal   & 1,714,390 \\
    Final usable dataset              & [XX] \\
    \bottomrule
  \end{tabular}
\end{table}

\subsection{Feature and Target Selection}
% Input features (X): pressure (hPa), CO2 (ppm), temperature (°C), humidity (%), PM2.5 (µg/m³)
% Target variable (y): exp_pl — expected path loss in dB
% This is the same target used in the prior project's XGBoost model

%% ── 4.4 ─────────────────────────────────────────────────────────────────────
\section{TinyML Model Architecture}
\label{sec:model-arch}

\subsection{Architecture Selection}
% Dense(5 → 8, ReLU) → Dense(8 → 1, Linear)
% 5 inputs: pressure (hPa), CO2 (ppm), temp (°C), humidity (%), PM2.5 (µg/m³)
% 1 output: predicted expected path loss (dB)
% Total parameters: (5×8 + 8) + (8×1 + 1) = 57
% Model size after int8 quantization: ~2.7 KB (model.tflite = 2704 bytes)

\subsection{Feature Normalization}
% StandardScaler: zero mean, unit variance
% Parameters exported from train_model.py, hardcoded in Arduino sketch:
% featureMeans = {1009.2605, 542.0334, 21.9533, 36.1344, 1.8987}
% featureStds  = {30.9770, 132.9057, 2.9012, 6.6737, 2.3336}
% Critical: device MUST normalize identically to training

\subsection{Path Loss to Link State Mapping}
% Model predicts continuous path loss (dB)
% Thresholds derived from exp_pl distribution analysis of prior dataset:
%   exp_pl < 117 dB → Good link
%   117 ≤ exp_pl < 133 dB → Degraded link
%   exp_pl ≥ 133 dB → Poor link
% These thresholds determine transmission behavior

%% ── 4.5 ─────────────────────────────────────────────────────────────────────
\section{Model Training Pipeline}
\label{sec:training-pipeline}

\subsection{Centralized Baseline Training}
% Architecture: Dense(5→8, ReLU) → Dense(8→1, Linear)
% Optimizer: Adam, learning_rate = 0.001
% Loss: Mean Squared Error
% Epochs: 50, batch size: 32
% Train/test split: 80/20, stratified by device_id
% StandardScaler normalization

\subsection{TFLite Conversion and Quantization}
% Post-training int8 quantization
% Representative dataset: 500 samples for calibration
% converter.target_spec.supported_ops = [TFLITE_BUILTINS_INT8]
% Keep float32 I/O for easier on-device use
% Result: model.tflite = 2,704 bytes

\subsection{C Header Generation}
% model.tflite → model.h (C byte array)
% alignas(8) const unsigned char g_model[]
% model.h = 17,723 bytes (hex representation)
% Normalization parameters exported as featureMeans[] and featureStds[]

%% ── 4.6 ─────────────────────────────────────────────────────────────────────
\section{Federated Learning Protocol}
\label{sec:fl-protocol}

\subsection{Communication Round Structure}
% FL round duration: 24 hours
% Minimum clients for aggregation: 3 of 6
% Each device independently decides when to send based on local timer
% FIGURE: FL round timeline showing local training, uplink, aggregation, downlink

\subsection{Weight Quantization and Payload Design}
% float32 → int8 quantization with scale factor 127
% Three message types:

\begin{table}[H]
  \centering
  \caption{LoRaWAN payload types and byte-level structure.}
  \label{tab:payloads}
  \begin{tabular}{llcp{5cm}}
    \toprule
    \textbf{Type} & \textbf{Code} & \textbf{Size} & \textbf{Content} \\
    \midrule
    Status report   & 0x00 & 8\,B  & link\_state, PDR, DR, packet\_count \\
    Model update    & 0x01 & 52\,B & version, num\_weights, 48 int8 deltas \\
    Global model    & 0x02 & 51\,B & version, num\_weights, 47 int8 weights \\
    \bottomrule
  \end{tabular}
\end{table}

% COMPARISON WITH TORRES SANCHEZ:
% They needed 1.39 KB per round (32 neurons AE) → 7–28 LoRaWAN messages
% We need 52 bytes per round (57 params, quantized) → 1 single LoRaWAN message
% This is possible because our model is ~25× smaller

\subsection{Event-Driven Transmission Strategy}
% Adaptive intervals based on predicted link state:
%   Good: 5 min intervals (NORMAL_TX_INTERVAL)
%   Degraded: 2 min intervals
%   Poor: 1 min intervals (URGENT_TX_INTERVAL)
% Event trigger: immediate TX when link transitions to POOR

%% ── 4.7 ─────────────────────────────────────────────────────────────────────
\section{Federated Learning Simulation}
\label{sec:fl-simulation}

\subsection{Simulation Environment}
% Following the methodology of Torres Sanchez et al., who state:
%   "The FL framework in this experiment simulates distributed devices
%    functioning as clients..."
% Our environment:
%   - Hardware: [CPU, RAM — fill in your PC specs]
%   - OS: Windows with Conda environment
%   - Python 3.13, TensorFlow [version], NumPy, scikit-learn
%   - Custom FedAvg implementation (not Flower — our model is small enough
%     that a lightweight custom implementation suffices)

\subsection{Non-IID Data Partitioning}
% Unlike random splits, we partition by device_id (ED0–ED5)
% Each virtual client receives the actual measurement data from its
% corresponding physical device location
% This produces naturally non-IID distributions because each room
% has different environmental characteristics

\begin{table}[H]
  \centering
  \caption{Data distribution per virtual client (non-IID split by device).}
  \label{tab:data-distribution}
  \begin{tabular}{lccc}
    \toprule
    \textbf{Client} & \textbf{Device} & \textbf{Location} & \textbf{Samples} \\
    \midrule
    Client 1 & ED0 & [Location] & [XX] \\
    Client 2 & ED1 & [Location] & [XX] \\
    Client 3 & ED2 & [Location] & [XX] \\
    Client 4 & ED3 & [Location] & [XX] \\
    Client 5 & ED4 & [Location] & [XX] \\
    Client 6 & ED5 & [Location] & [XX] \\
    \midrule
    \textbf{Total} & & & \textbf{[XX]} \\
    \bottomrule
  \end{tabular}
\end{table}

\subsection{FL Training Protocol}
% Each virtual client trains a local model on its partition for E local epochs
% After local training, weight updates aggregated via FedAvg
% Process repeats for N communication rounds
% Global model evaluated on held-out test set after each round
% Hyperparameters explored: rounds (5,10,20,50), local epochs (1,3,5),
%   min clients (2,3,4,6)

%% ── 4.8 ─────────────────────────────────────────────────────────────────────
\section{Edge Device Firmware}
\label{sec:firmware}

\subsection{TFLite Micro Integration}
% Model loaded from g_model[] in model.h
% AllOpsResolver (could optimize with MutableOpResolver)
% Tensor arena: 8192 bytes (8 KB)
% Static interpreter allocation

\subsection{Inference Pipeline}
% readSensors() → normalizeFeatures() → runInference() → threshold to link state
% Input: 5 float32 values → output: 1 float32 (predicted path loss dB)
% Link state derived: <117 Good, 117–133 Degraded, ≥133 Poor

\subsection{Local Training and Weight Management}
% Circular buffer: 32 samples (struct Sample with features[5] + label)
% Minimum 16 samples before training triggers
% 3 local epochs, simplified gradient computation
% Weight delta = localWeights - globalWeights
% Quantization: float → int8 (scale 127)

\subsection{LoRaWAN Communication}
% OTAA join with EU868
% FPort 3, ADR enabled
% Uplink: model update (52 bytes) — confirmed
% Downlink: global model (51 bytes) — via modem.available()

%% ── 4.9 ─────────────────────────────────────────────────────────────────────
\section{FL Aggregation Server}
\label{sec:server-design}

\subsection{Flask API Architecture}
% Endpoints: /health, /status, /uplink (TTN webhook), /downlink/<id>, /schedule_downlink/<id>
% Thread-safe FederatedLearningServer class with threading.Lock
% Flask application with 5 endpoints
% FederatedLearningServer class with thread-safe aggregation
% TTN webhook → base64 decode → dequantize → FedAvg → downlink preparation
% Round management: auto-start new round after aggregation

\subsection{FedAvg Implementation}
% Weighted average: w_{t+1} = Σ (n_k / Σ n_j) × w_k
% Currently uniform weighting (all clients report n_k = 1)

\subsection{TTN Webhook Integration}
% TTN v3 webhook → POST /uplink with JSON payload
% Base64-encoded frm_payload → decode → dequantize → aggregate
